% Generated by scripts/assemble.py; corrections belong in source/corrections.json.
\examhead{91}{南京师范大学 2024 年数学分析试卷}{89}
% p090-b21; source PDF page 90
\par\Needspace{5\baselineskip}
\blockanchor{p090-b21}
\majorquestion{一}计算题. 每题 5 分, 共 20 分.\par
% p090-b22; source PDF page 90
\blockanchor{p090-b22}
(1) 求极限  $\lim_{n\to\infty}\tan^{n}\left(\frac{\pi}{4}+\frac{1}{n}\right)$ .\par
% p090-b23; source PDF page 90
\blockanchor{p090-b23}
(2) 已知 $\lim_{x\to -\infty}(f(x) - f'(x)) = 2023$，求 $\lim_{x\to -\infty}f(x)$.\par
% p090-b24; source PDF page 90
\blockanchor{p090-b24}
(3) 求定积分 $\int_0^a x^2\sqrt{\frac{a - x}{a + x}}\mathrm{d}x (a > 0)$.\par
% p090-b25; source PDF page 90
\blockanchor{p090-b25}
(4) 交换积分顺序 \displayformula{\int_{\frac{1}{2}}^{\frac{1}{\sqrt{2}}} \mathrm{d}x \int_{\frac{1}{2}}^{x} f(x, y) \mathrm{d}y + \int_{\frac{1}{\sqrt{2}}}^{1} \mathrm{d}x \int_{x^2}^{x} f(x, y) \mathrm{d}y.}\par
% p091-b00; source PDF page 91
\par\Needspace{5\baselineskip}
\blockanchor{p091-b00}
\majorquestion{二}(15 分) 设级数 $\sum_{n=1}^{\infty} a_n$ 收敛于 $a$.\par
% p091-b01; source PDF page 91
\blockanchor{p091-b01}
(1) 求极限  $\lim_{n\to\infty}\frac{na_{1}+(n-1)a_{2}+\cdots+a_{n}}{n}.$\par
% p091-b02; source PDF page 91
\blockanchor{p091-b02}
(2) 证明: $\lim_{n\to \infty}\frac{a_1 + 2a_2 + \cdots + na_n}{n} = 0.$\par
% p091-b03; source PDF page 91
\par\Needspace{5\baselineskip}
\blockanchor{p091-b03}
\majorquestion{三}(15 分) 设 $f(x)$ 在 $x = 0$ 处连续, 且对一切 $x, y \in \mathbb{R}$, 有 $f(x + y) = f(x) + f(y)$. 证明: $f(x)$ 在 $\mathbb{R}$ 上连续且 $f(x) = f(1)x$.\par
% p091-b04; source PDF page 91
\par\Needspace{5\baselineskip}
\par\Needspace{6\baselineskip}
\blockanchor{p091-b04}
\majorquestion{四}(15 分) 设 $f(x)$ 在 $[a, b]$ 上二次可微, $f' \left( \frac{a + b}{2} \right) = 0$, 证明: 存在 $\xi \in (a, b)$, 使得\par
\nopagebreak[4]
% p091-b05; source PDF page 91
\blockanchor{p091-b05}
\displayformula{| f ^ {\prime \prime} (\xi) | \geq \frac {4}{(b - a) ^ {2}} | f (b) - f (a) |.}
% p091-b06; source PDF page 91
\par\Needspace{5\baselineskip}
\blockanchor{p091-b06}
\majorquestion{五}(15 分) 解答如下问题:\par
% p091-b07; source PDF page 91
\par\Needspace{6\baselineskip}
\blockanchor{p091-b07}
(1) 证明: 若 $f(x), g(x)$ 均在 $[a, b]$ 上可积, 则有\par
\nopagebreak[4]
% p091-b08; source PDF page 91
\blockanchor{p091-b08}
\displayformula{\left(\int_ {a} ^ {b} f (x) g (x) \mathrm{d} x\right) ^ {2} \leq \int_ {a} ^ {b} f ^ {2} (x) \mathrm{d} x \int_ {a} ^ {b} g ^ {2} (x) \mathrm{d} x.}
% p091-b09; source PDF page 91
\par\Needspace{6\baselineskip}
\blockanchor{p091-b09}
(2) 证明: 若 $f'(x)$ 在 $[0,1]$ 上连续, $f(0) = 0$, 则有\par
\nopagebreak[4]
% p091-b10; source PDF page 91
\blockanchor{p091-b10}
\displayformula{\int_ {0} ^ {1} f ^ {2} (x) \mathrm{d} x \leq \frac {1}{2} \int_ {0} ^ {1} [ f ^ {\prime} (x) ] ^ {2} \mathrm{d} x.}
% p091-b11; source PDF page 91
\par\Needspace{5\baselineskip}
\blockanchor{p091-b11}
\majorquestion{六}(10 分) 设函数 $f(x)$ 在 $[a, +\infty)$ 上单调递减, 且 $f(x) > 0, x \in [a, +\infty)$, 证明: 反常积分 $\int_{a}^{+\infty} f(x) \mathrm{d}x$ 与 $\int_{a}^{+\infty} f(x) \sin^2 x \mathrm{d}x$ 敛散性相同.\par
% p091-b12; source PDF page 91
\par\Needspace{5\baselineskip}
\blockanchor{p091-b12}
\majorquestion{七}(20 分) 解答如下问题:\par
% p091-b13; source PDF page 91
\par\Needspace{6\baselineskip}
\blockanchor{p091-b13}
(1) 已知 $f(x, y)$ 在点 $(0, 0)$ 的某邻域内连续, 且\par
\nopagebreak[4]
% p091-b14; source PDF page 91
\blockanchor{p091-b14}
\displayformula{\lim _ {(x, y) \rightarrow (0, 0)} \frac {f (x , y) - x y ^ {2}}{1 - \cos \sqrt {x ^ {2} + y ^ {2}}} = 1.}
% p091-b15; source PDF page 91
\blockanchor{p091-b15}
试问：(0,0) 是否为 $f(x,y)$ 的驻点(或稳定点)? 进一步讨论 (0,0) 是否为 $f(x,y)$ 的极值点? 若是, 是极大值点还是极小值点? 若不是, 请说明理由.\par
% p091-b16; source PDF page 91
\blockanchor{p091-b16}
(2) 设 $u = u(x, y, z)$, $x = r \sin \varphi \cos \theta$, $y = r \sin \varphi \sin \theta$, $z = r \cos \varphi$, 若 $x u_{x} + y u_{y} + z u_{z} = 0$, 证明 $u$ 仅为 $\theta$ 与 $\varphi$ 的函数.\par
% p091-b17; source PDF page 91
\par\Needspace{5\baselineskip}
\blockanchor{p091-b17}
\majorquestion{八}(15 分) 讨论级数 $\sum_{n=1}^{\infty} \frac{(-\alpha)^n n!}{n^n} (\alpha > 0)$ 的敛散性, 若收敛, 请指出是条件收敛还是绝对收敛.\par
% p091-b18; source PDF page 91
\par\Needspace{5\baselineskip}
\par\Needspace{6\baselineskip}
\blockanchor{p091-b18}
\majorquestion{九}(15 分) 计算曲面积分\par
\nopagebreak[4]
% p091-b19; source PDF page 91
\blockanchor{p091-b19}
\displayformula{\iint_ {S} \left(x ^ {2} \cos \alpha + y ^ {2} \cos \beta + z ^ {2} \cos \gamma\right) \mathrm{d} S.}
% p091-b20; source PDF page 91
\blockanchor{p091-b20}
其中 $S$ 为圆锥曲面 $x^{2} + y^{2} = z^{2}$ 上位于 $0 \leq z \leq h$ 的一部分, 其法向量恒与 $z$ 轴正向相交成锐角, $\cos \alpha, \cos \beta, \cos \gamma$ 为法向量的方向余弦.\par
% p091-b21; source PDF page 91
\par\Needspace{5\baselineskip}
\blockanchor{p091-b21}
\majorquestion{十}(10 分) 证明: $\int_{0}^{+\infty} \frac{\cos x^{2}}{x^{y}} \mathrm{d}x$ 关于参变量 $y$ 在 $(-1,1)$ 内内闭一致收敛.\par
